% GDZ PPN236018647,native568lower–571upper;HKUST569–572.
\providecommand{\DiaryBracketNote}[2]{\textnormal{[#1)]}\begingroup\renewcommand{\thefootnote}{[#1)}\stepcounter{footnote}\footnotetext{#2]}\endgroup}
\providecommand{\DiaryMoon}{\tikz[baseline=-.15ex,x=.65ex,y=.65ex]{\draw[line width=.35pt] (.32,.6)..controls(-.42,.6)and(-.42,-.6)..(.32,-.6)..controls(-.18,-.44)and(-.18,.44)..(.32,.6);}}
\providecommand{\DiaryAquarius}{\tikz[baseline=-.1ex,x=.4ex,y=.4ex]{\draw[line width=.45pt](0,.7)--(.3,.9)--(.6,.7)--(.9,.9)--(1.2,.7)--(1.5,.9);\draw[line width=.45pt](0,.1)--(.3,.3)--(.6,.1)--(.9,.3)--(1.2,.1)--(1.5,.3);}}
\NumberedItem{97}
\indent Formulas novas \textls[120]{exactas} pro parallaxi eruimus.
\par\noindent\hfill Br[unsvigae, 1799] Apr. 8.

\medskip\begingroup\small
\indent Diese Formeln (für die Mondparallaxe) hat \textsc{Gauss} an verschiedenen Stellen aufgezeichnet, und zwar wie es seine Gewohnheit war, auf leeren Blättern von Büchern seiner Bibliothek. In unmittelbarem Zusammenhang mit unserer Tagebuchnotiz scheint die folgende Eintragung zu stehen, die sich auf den hintern Einbandseiten von \textsc{Klügels} \textit{Analytischer Geometrie}, 1770, (Nr. 739 der \textsc{Gauss}bibliothek) findet:
\par\begingroup\normalsize
\indent Ist des Mondes Declination \(=d\), sein Stundenwinkel \(=h\), Horizontalparallaxe \(=\pi\), Polhöhe \(=\varphi\), so ist (genau)
\[
\operatorname{cotg}\text{ par[allaxis] Asc[ensionis] r[ectae]}=\frac{\cos d}{\cos\varphi\sin h\sin\pi}-\operatorname{cotg}h
\]
\indent Diameter im Horizont \(=\Delta\)
\par\indent Vergrösserter Diameter \(=\Delta'\)
\[
\begin{aligned}
\sin\tfrac12\Delta'&=\frac{\sin\frac12\Delta\sin(h+p)\cos d'}{\sin h\cos d}\\[1ex]
\operatorname{tang}d'&=\frac{\sin(h+p)}{\sin h}\left(\operatorname{tang}d-\frac{\sin\pi\sin\varphi}{\cos d}\right)
\end{aligned}
\]
\begin{center}\fbox{\strut REV. GALEN}\ \DiaryBracketNote{*}{Siehe die unerklärliche Buchstabenverbindung GALEN auch oben S. 187 aus Scheda Ac.}\end{center}
\par\endgroup

Am ausführlichsten ist von diesen Formeln in einem Briefe die Rede, den \textsc{Gauss} in Braunschweig am 17. April 1799 geschrieben hat und der gegenwärtig in der Autographensammlung \textsc{Darmstaedter} der Königlichen Bibliothek zu Berlin aufbewahrt wird. Es heißt in diesem Briefe:
\par\bigskip\noindent\hfill68*\hspace*{.14\linewidth}
\par\endgroup

\clearpage\SourceLeaf{569}\ArithmeticHead{540}{L}{TAGEBUCH MIT ERLÄUTERUNGEN.}
\indent Ich habe seit meinem letzten Briefe das Brouillon wieder gefunden von der Berechnung der Bedeckung des \(\tau\) \DiaryAquarius\DiaryBracketNote{*}{\(\tau\) aquarii.}. Ich hatte diese meistens ad normam \textsc{Bohnenberger}\DiaryBracketNote{**}{\textit{Anleitung zur geographischen Ortsbestimmung vorzüglich vermittelst des Spiegelsextanten} von M. J. G. F. \textsc{Bohnenberger}, Göttingen 1795, §§ 187—190, S. 341 ff.} geführt, weil ich damals mich nicht gleich in diesen Arbeiten genug orientirt hatte, um Ihnen das Resultat sogleich angeben zu können, wie Sie damals wünschten. Daraus sehe ich also, wie \textsc{Bohnenberger} die Längenparallaxe u.s.w. berechnet hat (mit Ausnahme der sphäroidischen Gestalt der Erde, auf die ich damals nicht Rücksicht genommen hatte, weil ich irrig voraussetzte, die Horizontalpar[allaxe] in dem Berliner Jahr B[uch] gelte für Berlin) und dass seine Methode nicht Näherung sei. In der That sind auch die exacten Formeln für Jemand, der in der analytischen Trigon[ometrie] nur mässig geübt ist, so leicht zu finden, dass es unbegreiflich sein würde, wenn sie noch von Niemand entwickelt wären. Indess muss ich gestehen, dass die \textls[120]{Form}, in der B[\textsc{ohnenberger}] sie darstellt\DiaryBracketNote{***}{Siehe a. a. O. S. 346 unten, beziehungsweise S. 350.}, mir weniger bequem scheint als die meinige, obgleich er vielleicht das Gegentheil geglaubt haben mag. Er richtet es nemlich so ein, dass man alles vermittelst der trigonometrischen Tafeln machen kann; allein ich glaube, dass der Grund davon ein Vorurtheil ist. .\ .\ . Meine Formeln werde ich also zwar nicht bekannt machen, zumal da sie wahrscheinlich schon in andern Büchern stehen werden (vielleicht in \textsc{Cagnoli}\DiaryBracketNote{\ensuremath{\dagger}}{\textit{Traité de Trigonométrie rectiligne et sphérique etc.} par M. \textsc{Cagnoli}, traduit de l’Italien par M. \textsc{Chompré}, Paris 1786.}) aber rechnen werde ich doch immer nach ihnen. Dies sind sie:
\par\smallskip
\begin{tabular}{@{\hspace{\parindent}}ll@{}}
\(\pi\)&Horizontalparallaxe unterm Aequator\\
\(\pi'\)&verbesserte Horiz. Par.\\
\(\rho\)&Entfernung vom Mittelpunkte der Erde\\
\(\beta\)&Breite des Zeniths\\
\(\lambda\)&Länge des Zeniths\\
\(b\)&Breite des \DiaryMoon\ [Mondes]\\
\(l\)&Länge des \DiaryMoon\\
\multicolumn{2}{@{\hspace{\parindent}}l@{}}{\(\delta=l-\lambda\)}\\
\(r\)&Halbmesser des \DiaryMoon
\end{tabular}

\clearpage\SourceLeaf{570}\ArithmeticHead{541}{R}{1799 APR. 8.}
\noindent\begin{tabular}{@{\hspace{\parindent}}ll@{}}
\(r'\)&Vergrösserter Halbmesser\\
\(b'\)&Scheinbare Breite [des Mondes]\\
\(p\)&Längenparallaxe
\end{tabular}
\[
\begin{array}{r@{\quad}rcl}
\mathrm{I.}&\sin\pi'&=&\rho\sin\pi\\[1ex]
\mathrm{II.}&\operatorname{cotang}p&=&\displaystyle\frac{\cos b}{\cos\beta\sin\delta\sin\pi'}-\operatorname{cotang}\delta\\[2ex]
\mathrm{III.}&\operatorname{tang}b'&=&\displaystyle\frac{\sin\delta+p}{\sin\delta}\left(\operatorname{tang}b-\frac{\sin\pi'\sin\beta}{\cos b}\right)\\[2ex]
\mathrm{IV.}&\sin r'&=&\displaystyle\frac{\sin\delta+p}{\sin\delta}\,.\,\frac{\cos b'}{\cos b}\sin r\text{[*)].}
\end{array}
\]
\begingroup\renewcommand{\thefootnote}{[*)}\stepcounter{footnote}\footnotetext{Die Formeln II, III, IV stimmen für \(\rho=1,\ \beta=\varphi,\ b=d,\ \delta=h,\ r=\frac12\Delta,\ r'=\frac12\Delta',\ b'=d'\) mit den oben aus \textsc{Klügel} mitgeteilten überein. \textsc{Gauss} hat sie mit den nachfolgenden Gebrauchsformeln in den hier benutzten Bezeichnungen in seinen 1798 erworbenen Abdruck von \textsc{Röslers} \textit{Handbuch der praktischen Astronomie}, 1788, eingeschrieben.]}\endgroup

Zur Ausübung sind in den meisten Fällen folgende Formeln die bequemsten.

I. Man suche den Logarithmen einer Zahl \(M\) nach der Formel
\[
\frac1M=\cos b-\rho\sin\pi\cos\beta\cos\delta,
\]
\noindent dann wird
\[
\begin{array}{r@{\quad}rcl}
\mathrm{II.}&\operatorname{tang}p&=&M\rho\sin\pi\cos\beta\sin\delta\\
\mathrm{III.}&\operatorname{tang}b'&=&M\sin b\cos p-M\rho\cos p\sin\beta\sin\pi\\
\mathrm{IV.}&\sin r'&=&M\cos b'\cos p\sin r,
\end{array}
\]
\begin{quote}
(wofür man \textls[120]{allemal} die Näherung
\[
r'=M\cos b'\cos p\,.\,r
\]
\noindent brauchen darf).
\end{quote}
Nach den letztern Formeln habe ich die Berechnung für Seeberg als Beispiel beigelegt. Ich merke nur noch an, dass das was \textsc{Bohnenberger} \(N\) nennt (nach meinem Brouillon\DiaryBracketNote{**}{Die Bezeichnung \(N\) kommt bei \textsc{Bohnenberger} nicht vor; \textsc{Gauss} hatte wohl die auf S. 350 der \textit{Anleitung} stehenden Näherungsformeln vor Augen, in denen der gemeinsame Nenner \(\cos B\,(\sin\frac12 A)^2\) in der Bezeichnung von \textsc{Gauss} gleich \(\frac1{2M}\) ist.}) immer gleich ist dem Logarithm von \(2M\) .\ .\ .\ .\ .

\medskip\begingroup\small
\renewcommand{\thefootnote}{***)}
\indent Die \textsc{Gauss}schen Formeln finden sich im Wesentlichen bei \textsc{Bohnenberger} a. a. O. § 187, S. 344—346\footnote{\textsc{Gauss}’ Bezeichnungen weichen von denen \textsc{Bohnenbergers} etwas ab, vergl. die Zusammenstellung der Bezeichnungen bei \textsc{Bohnenberger} a. a. O. S. 345.}; die Formel für die Längenparallaxe (II. der \textsc{Gauss}schen Formeln) hat nach \textsc{Bohnenbergers} Zitat bereits »\textsc{Lexell} in den Berliner Ephemeriden für 1777, S. 152 u. f.« bekannt gemacht.
\par\endgroup

\clearpage\SourceLeaf{571}\ArithmeticHead{542}{L}{TAGEBUCH MIT ERLÄUTERUNGEN.}
\begingroup\small
\indent Die zu Anfang unserer Briefstelle erwähnte Bedeckung von \(\tau\) aquarii ist wohl diejenige, die nach den Allgemeinen geographischen Ephemeriden \textbf{3}, 1799, S. 202 am 13. Dezember 1798 stattgefunden hat und von dem damaligen K. Preuß. Obersten und General-Quartiermeister bei der Neutralitäts-Armee \textsc{von Lecoq} in Preußisch-Minden beobachtet worden ist. Die Vergleichung mit den im \textsc{Gauss}archiv befindlichen Briefen \textsc{von Lecoqs} an \textsc{Gauss} zeigt, daß der hier im Auszug mitgeteilte Brief vom 17. April 1799 an \textsc{von Lecoq} gerichtet ist. In seinem Bericht \textit{Über die trigonometrische Aufnahme in Westphalen}\footnote{Monatliche Correspondenz \textbf{8}, 1803, S. 136 ff., S. 197 ff.} schreibt \textsc{von Lecoq} (S. 139): »Im astronomischen Theile ist mir der Doctor \textsc{Gauss} von großem Nutzen gewesen; seine Ausrechnungen und Briefe haben zu meinem Unterrichte viel beygetragen .\ .\ .«.
\par\noindent\hfill\textsc{Galle.}
\par\endgroup
